% UC10 --- Visualizar Notificações

\begin{longtable}{| p{.40\textwidth} | p{.50\textwidth} |}
\hline

\textbf{Nome do Caso de Uso} &
UC10 --- Visualizar Notificações \\
\hline

\textbf{Ator Principal} &
Usuário \\
\hline

\textbf{Atores Secundários} &
--- \\
\hline

\textbf{Resumo} &
Permite ao usuário consultar os lembretes recebidos pelo sistema, como
lembretes diários de registro de humor, lembretes de medicamento e alertas de
indicadores. \\
\hline

\textbf{Pré-condições} &
Usuário autenticado. \\
\hline

\textbf{Pós-condições} &
O usuário visualiza as notificações e pode ser direcionado à tela
correspondente. \\
\hline

\multicolumn{2}{|c|}{\textbf{Cenário Principal}}\\
\hline

\textbf{Ações do Ator} &
\textbf{Resposta do Sistema} \\
\hline

1. Acessa a tela de notificações. & \\
\cline{1-2}

&
2. O sistema apresenta as notificações do usuário, das mais recentes para as
mais antigas. \\
\cline{1-2}

3. Toca em uma notificação. & \\
\cline{1-2}

&
4. O sistema encaminha para a tela correspondente à notificação. \\
\hline

\multirow{2}{.40\textwidth}{\textbf{Restrições\\e validações}}
&
Somente as notificações do usuário autenticado são exibidas. \\
\cline{2-2}

&
Notificações fora do período de validade são exibidas como expiradas. \\
\hline

\multicolumn{2}{|c|}{\textbf{Cenário Alternativo A1 --- Notificação recebida via push}}\\
\hline

\textbf{Ações do Ator} &
\textbf{Resposta do Sistema} \\
\hline

1. Toca na notificação recebida no dispositivo. & \\
\cline{1-2}

&
2. O sistema abre a tela de origem da notificação. \\
\hline

\multicolumn{2}{|c|}{\textbf{Cenário de Exceção E1 --- Nenhuma notificação}}\\
\hline

\textbf{Ações do Ator} &
\textbf{Resposta do Sistema} \\
\hline

1. Acessa a tela de notificações. & \\
\cline{1-2}

&
2. O sistema informa que não há notificações. \\
\hline

\end{longtable}
